\begin{lemma}[A symmetric bilinear form is controlled by its diagonal]
Let $E$ be a finite-dimensional real inner-product space, let $S$ be a continuous
bilinear form on $E$ that is symmetric, $S(a,b)=S(b,a)$ for all $a,b\in E$, and
let $C\ge0$ satisfy
\[
  \bigl|S(a,a)\bigr|\le C|a|^2\qquad(a\in E).
\]
Then $\|S\|\le C$, where $\|S\|$ is the operator norm, i.e.\ the least constant
$K$ with $|S(a,b)|\le K|a|\,|b|$ for all $a,b\in E$.
\end{lemma}
\begin{proof}
It suffices to prove
\begin{equation}\label{eq:sfnd-goal}
  \bigl|S(a,b)\bigr|\le C|a|\,|b|\qquad(a,b\in E),
\end{equation}
since $C\ge0$ and the operator norm of a continuous bilinear form is the least
such constant.

\emph{Step 1: the polarization bound.}
Since $S$ is bilinear and symmetric,
\[
  S(u+v,u+v)-S(u-v,u-v)
   =\bigl(S(u,u)+2S(u,v)+S(v,v)\bigr)
     -\bigl(S(u,u)-2S(u,v)+S(v,v)\bigr)=4S(u,v),
\]
the cross terms being combined by $S(v,u)=S(u,v)$. Hence
\[
  S(u,v)=\tfrac14\bigl(S(u+v,u+v)-S(u-v,u-v)\bigr),
\]
and the diagonal hypothesis applied at $u+v$ and at $u-v$ gives
\[
  \bigl|S(u,v)\bigr|
   \le\tfrac14\Bigl(C|u+v|^2+C|u-v|^2\Bigr)
   =\tfrac14\,C\cdot2\bigl(|u|^2+|v|^2\bigr)
   =\frac{C\bigl(|u|^2+|v|^2\bigr)}{2},
\]
using the parallelogram law $|u+v|^2+|u-v|^2=2(|u|^2+|v|^2)$, which follows by
adding the two expansions $|u\pm v|^2=|u|^2\pm2\langle u,v\rangle+|v|^2$. So
\begin{equation}\label{eq:sfnd-half}
  \bigl|S(u,v)\bigr|\le\frac{C\bigl(|u|^2+|v|^2\bigr)}{2}\qquad(u,v\in E).
\end{equation}

\emph{Step 2: the degenerate cases.}
If $a=0$ then $S(a,b)=S(0,b)=0$ by linearity of $S$ in its first argument, so
$|S(a,b)|=0\le C|a|\,|b|$, the right-hand side being $\ge0$ because $C\ge0$.
Likewise if $b=0$. So assume $a\ne0$ and $b\ne0$, hence $|a|>0$ and $|b|>0$.

\emph{Step 3: scaling the two arguments against each other.}
Put
\[
  t=\sqrt{\frac{|b|}{|a|}}>0,\qquad\text{so}\qquad
  t^2=\frac{|b|}{|a|},\qquad t^{-2}=\frac{|a|}{|b|} .
\]
By bilinearity, $S(ta,t^{-1}b)=t\,t^{-1}S(a,b)=S(a,b)$. Applying
\eqref{eq:sfnd-half} to $u=ta$ and $v=t^{-1}b$, and using
$|ta|=t|a|$ and $|t^{-1}b|=t^{-1}|b|$ (both scalars being positive),
\[
  \bigl|S(a,b)\bigr|
   =\bigl|S(ta,t^{-1}b)\bigr|
   \le\frac{C\bigl(t^2|a|^2+t^{-2}|b|^2\bigr)}{2}
   =\frac{C\bigl(\tfrac{|b|}{|a|}|a|^2+\tfrac{|a|}{|b|}|b|^2\bigr)}{2}
   =\frac{C\bigl(|a|\,|b|+|a|\,|b|\bigr)}{2}
   =C|a|\,|b| ,
\]
which is \eqref{eq:sfnd-goal}. (The scaling is exactly what replaces the
arithmetic mean $(|u|^2+|v|^2)/2$ of Step 1 by the geometric mean
$|a|\,|b|$; without it \eqref{eq:sfnd-half} alone gives only the weaker bound
$|S(a,b)|\le C$ on the unit ball.)
\end{proof}
